\documentclass[12pt,letterpaper]{article}

\usepackage[margin=0.75in]{geometry}
\usepackage{amsmath,amssymb}
\usepackage{enumitem}

\pagestyle{empty}
\setlength{\parindent}{0pt}
\setlength{\parskip}{0.5em}
\setlength{\emergencystretch}{2em}

\begin{document}

\begin{center}
  {\large\bfseries MATH 15200: Calculus II}\\
  {\large\bfseries Homework 5: Inverse Functions and Techniques of Integration}
\end{center}

\vspace{0.25em}
\hrule
\vspace{0.8em}

\textbf{Due:} Monday, November 16, 2026, at 11:59 PM.

This is the written portion of Homework 5; complete the routine online portion
separately in Edfinity. Upload your written solutions to Gradescope, via Canvas. You may
collaborate, but you must write and submit your own solutions independently.



\vspace{0.8em}


\begin{enumerate}[label=\textbf{\arabic*.},leftmargin=2em,itemsep=1.25em]
  \item Let $f(x)=x+e^x$. Since $f'(x)=1+e^x>0$, the function $f$ is
    increasing, so it has an inverse function $f^{-1}$. There is no formula
    for $f^{-1}$ in terms of familiar functions. Nevertheless, evaluate
    \[
      \int_1^{1+e}f^{-1}(y)\,dy.
    \]

    \emph{Hint:} First show that 
    \[
      \int_{f(a)}^{f(b)}f^{-1}(y)\,dy=b\,f(b)-a\,f(a)-\int_a^b f(x)\,dx.
    \]
    To do this, make the substitution $y=f(x)$ on the left side, and then
    integrate by parts. Then find $a$ and $b$ with $f(a)=1$ and $f(b)=1+e$.

  \item For all $x \neq 0$, set 
    \[
      g(x)=\arctan x+\arctan\Bigl(\frac1x\Bigr).
    \]
    Show that $g'(x)=0$ for every $x\ne0$. Then show that
    $g(x)=\frac\pi2$ for $x>0$ and $g(x)=-\frac\pi2$ for $x<0$. Since $g'$
    is zero wherever $g$ is defined, why doesn't this contradict the fact that
    a function whose derivative is zero at every point of an interval is
    constant on that interval?

    \emph{Hint:} Once you know that $g$ is constant on $(0,\infty)$ and
    constant on $(-\infty,0)$, find each constant by evaluating $g$ at a
    convenient point of that interval.

  \item Let $a,b>0$. Find the area enclosed by the ellipse
    \[
      \frac{x^2}{a^2}+\frac{y^2}{b^2}=1.
    \]
    Check that when $a=b=r$, your answer is the area of a circle of radius
    $r$.

    \emph{Hint:} By symmetry, the area is four times the area of the part of
    the ellipse in the first quadrant. Solve for $y$ on the upper half of the
    ellipse, set up the integral, and use the substitution $x=a\sin\theta$.
\end{enumerate}

\end{document}
